\subsection{Held-Out Gate and Partial Rollback}
\label{app:procedure-validation}
After each full training pass, we evaluate the candidate library on $T_v$. For skill $s$, the gate uses the retrieved-support set $T_{v,s}$ defined in the main text and compares the candidate with the fixed cold-start predicate $p_{si}$. With one predicate per skill, the implemented validation score is the mean correctness of its Boolean judgments, equal to $1-R_v$.

As in the main formulation, the held-out gate uses the following strict admission rule:
\[
\begin{aligned}
&g_{\mathrm{strict}}(T_v,m_p^*,p_{si})\\
&\quad=\mathds{1}[R_v(m_p^*\mid s)<R_v(p_{si}\mid s)].
\end{aligned}
\]
Validation ties are rejected. Let $p_s^{(i)}$ denote the predicate at the start of the iteration, and let $\widetilde p_s$ denote its value after the training pass. For changed skills, the retained update is
\[
p_s^{(i+1)}=\begin{cases}
\widetilde p_s,&g_{\mathrm{strict}}(T_v,\widetilde p_s,p_{si})=1,\\
p_s^{(i)},&\text{otherwise}.
\end{cases}
\]
Skills without measurable validation support are also reverted. The rollback is per skill, so the retained library can contain both accepted replacements and iteration-start predicates. Comparing to $p_{si}$ does not ensure monotonic validation performance relative to $p_s^{(i)}$.

Training data and replay determine predicate proposals while validation determines which updates are retained. Test data is used only for evaluation, not for predicate repair or checkpoint selection. We distinguish scores measured before validation rollback from those of the retained library after rejected updates are reverted.